\section{Novelty Statement}

Unlike conventional GPU collision-detection methods that rely primarily on compute-shader hashing, sorting, atomic operations, or iterative neighbor-list construction, gPCD uses rasterization-stage parallelism to generate the potentially colliding set (PCS) directly in GPU memory. The resulting cell array is then consumed by a compute-shader narrow phase with one thread assigned to each particle.

A second novel feature is the explicit dual duplicate-elimination strategy. The method resolves both

\begin{enumerate}
    \item corner duplicates caused by multiple corners of the same particle mapping to the same cell, and
    \item particle duplicates caused by particle pairs sharing multiple cells.
\end{enumerate}

These duplicates are removed locally within each compute invocation using a per-particle collision cache, thereby reducing redundant distance evaluations without requiring global synchronization.

Additional novelty arises from the geometric occupancy constraint

\[
d_p \leq 1
\]

which guarantees that a particle may intersect no more than eight cells. This creates a bounded-memory neighbor-search structure with predictable access patterns and fixed per-particle occupancy complexity.

Collectively, these features form a hybrid rasterization--compute collision framework that differs from traditional spatial hashing and uniform-grid methods by leveraging graphics hardware scheduling for broad-phase candidate generation and compute parallelism for exact collision resolution.

\subsection{Summary}

gPCD combines rasterization-stage occupancy generation with compute-stage exact collision testing in a fully GPU-resident workflow. The graphics pipeline simultaneously performs point-particle rendering and broad-phase occupancy construction, while the compute stage performs exact collision evaluation.
